\subsection{Scope of the Preliminary Results}
\label{sec:results_model_presentation}

The results in this chapter are provisional development results included to
document the current state of the study. The detailed shower-separation plots
use separately trained \(2e\) and \(3e\) \texttt{ECalTpadTransformer} models.
Each model used a sample of \(20{,}000\) events, divided into \(16{,}000\)
training events, \(3{,}000\) validation events and \(1{,}000\) held-out test
events. The plots below are calculated from the validation split using the
checkpoint with the lowest validation loss. They should therefore be treated
as exploratory diagnostics rather than final, unbiased estimates of model
performance. The final comparison will use the held-out test samples after the
model choices have been fixed.

The training history and trigger-pad ablation were produced by an earlier,
similarly configured \(3e\) development run. They illustrate training and
ablation behaviour, but they are not provenance-matched to the checkpoints
used for the remaining validation figures.

For most shower-separation diagnostics, the numerical group labels are treated
as interchangeable. One whole-event permutation is chosen to maximise the
number of correctly grouped ECal hits, and that same permutation is then kept
for all layers, energy weights and confidence selections within the event.
The reported quantity is consequently a permutation-aligned grouping accuracy.
This evaluates whether the showers have been separated independently of the
names assigned to the groups. The training curves and trigger-pad ablation are
exceptions and use the native class labels of their respective run.

The currently available detailed figures focus on the Transformer with ECal
and trigger-pad-track inputs. The controlled comparison with the GravNet
models, the ECal-only variants and the exploratory multi-task model remains to
be added.




\clearpage
\subsection{Pending Comparisons}

The current figures establish the evaluation framework and show the behaviour
of the preliminary Transformer models, but they do not yet answer every
research question. The final chapter will add held-out test comparisons for
all four baseline models, including model size and inference-time information,
and will present the outputs of the exploratory multi-task model. No final
ranking of GravNet and Transformer architectures is inferred from the figures
presented here. Qualitatively similar trends have been observed during larger
ongoing runs, but their numerical results are not included until the complete
training and evaluation procedure has been fixed.
